\begin{scanpage}{0043}{26}
\sectiontitle{4}{The Field Radiated by the Nonuniform Part of the Current}
In \S\ 1 and 3 we represented the rigorous and approximate expressions for the diffraction field by integrals along the same contour in the complex variable plane. By subtracting the approximate expression from the rigorous expression, we find the field created by the nonuniform part of the current. It is determined by integrals of the type
\[\int_{D_0}p(\alpha,\phi,\phi_0,\zeta)e^{ikr\cos\zeta}\,d\zeta,\tag{4.01}\]
which, with the replacement of the variable $\zeta$ by $s=\sqrt2e^{i\frac\pi4}\sin\frac\zeta2$, are transformed to the form
\[e^{ikr}\int_{-\infty}^{\infty}q(\alpha,\phi,\phi_0,s)e^{-krs^2}\,ds\tag{4.02}\]
and may be approximately calculated by the method of steepest descents.

For this purpose, let us expand the function $q(s)$ into a Taylor series
\[q(\alpha,\phi,\phi_0,s)=q_0+q_1\cdot s+q_2\cdot s^2+\cdots\tag{4.03}\]
Let us note that expansion (4.03) does not have meaning only in the particular case
\[\left.\begin{aligned}\phi&=\pi\pm\phi_0&&\text{with }\phi_0=0;\ \pi,\\\phi&=2\alpha-\pi-\phi_0&&\text{with }\phi_0=\alpha-\pi,\end{aligned}\right\}\tag{4.04}\]
when the observation direction $(\phi)$ coincides with the direction of propagation of the incident wave glancing along one of the wedge faces.
\end{scanpage}
